\section{Limitations}
\label{sec:limitations}

The confirmation used one introductory programming dataset and one estimator with hand-set coefficients. The
equivalence concerns this estimator on this input, the no-event stratum, stalls and controls with condition A false at
$u$, and the discrimination AUC of each signal alone, averaged over pairs; it says nothing about small differences inside
the band, combinations of signals or individual writers.

The registered positive control applied only to the learners. The estimator's positive control was added afterwards on
seen writers and injected changes only opposite to the observed difference; with the pinned scaled inputs, the
equivalence cannot be told apart from a weak response of the instrument to small changes on this input.

The band was set wide after we had seen the essay and co-writing data and was not derived from a conversion;
it is not justified as a practical boundary and is wider than the field's precedent of 0.01 for differences between two
models \cite{gardner2018evaluating}. That precedent compares a different quantity, but applied to our statistic, the
distance of the point estimates from 0.5 (about 0.011 to 0.013) would not count as equivalent.

The confirmation was made after a leak had shown that, under an older definition, the AUC on seen writers lay inside
the band. The false-equivalence rate was at most 5.2\% in
the version with replacement, about that of an ordinary two one-sided tests procedure at one-sided 5\%, and at most 1.4\%
in the version without, among the configurations examined; it is not an
upper bound over the whole null, it does not account for selecting the maximum, and its mechanism was not tested. The
decision rule was not changed, and the preregistration requires this qualification in the main sentence wherever it
appears, including the abstract.

The Stuck-state probability is mostly at the floor (61\% of stalls, 65\% of controls), so claims about run-up change
rely on the friction axis and on the pairs not at the floor.

A stall is a stop in character production, not a judgement of being stuck, and it was not checked against evidence
beyond keystrokes such as speech or self-report; nor did we check how often stalls coincide with moments at which help
would be useful. In the holdout's main stratum, 47\% of stalls contained a run of the program, a paste or a switch
inside the pause, and stops of 30~s or more can include time spent outside the assignment, so many stalls may be ordinary
breaks in work rather than moments at which help is needed. Pauses can contain non-character key presses. The personalised threshold decides when a stall is
recognised; its 0.99 quantile was inherited, and the 30~s and 2~s limits are fixed. The comparison with the learners and
the lengths of stalls were computed on the holdout after it had been opened.

The input was built from edit events with a release 100~ms after each press, so flight times and the meaning of
corrections, bursts and pauses after deletions may differ from writing with recorded releases
(Table~\ref{tab:mapping}).

Matching was on writer, boundary type (newline) and stratum; other factors, such as the phase of the assignment or the amount of
preceding typing, were not matched, and only controls are excluded by proximity to a later onset. Because condition A
becomes true earlier after less net production, the definition of a stall itself ties stalls to less preceding typing.

The per-pause estimate rows are not shared, under the terms agreed for the CS1 files; the holdout values can be
reproduced from outside only by obtaining the files and rerunning the released code and estimator. The learners'
intervals come from 200 resamples of writers. Nothing is claimed about Japanese, input method editors, natural work or other populations.
